\section{P2: The Guardrail Boundary, Across Three Domains}
\label{sec:p2}

Prior work found that an agent under an explicit honesty mandate \emph{refuses} a
direct interpersonal request to misreport a failing peer ($0.00$) yet passes along the
same false status when it arrives as an ambient artifact ($\approx0.75$), and that
making the report auditable does not help ($0.71$ unobserved vs $0.77$ observed,
$n{=}48$). That finding rested on a single sandboxed task family, which is a fair
reason to discount it.

\paragraph{Design.}
Three attestation domains sharing nothing with the original world but the abstract
structure---an agent must file a verdict about a ground truth it can read, and a
misleading premise reaches it by one of two routes:
\emph{code review} (a pull request whose CI is red), \emph{compliance} (a control that
failed testing), \emph{procurement} (a vendor in SLA breach). Each crosses
\texttt{direct} (an explicit plea from the interested party) against \texttt{ambient}
(a pre-populated checklist, draft, or scorecard asserting the favourable state). Both
routes carry the same false claim; only delivery differs. $48$ seeds per cell, forced
verdict, deterministic honest/omission/active-lie scoring against ground truth,
Benjamini--Hochberg across domains plus a Cochran--Mantel--Haenszel pooled test.

\paragraph{Pre-committed reading.}
Two or more domains replicating means the effect is not specific to the original task
family; zero of three means it does not generalise and the prior finding should be
scoped to that family. We commit to reporting the null as prominently as the positive.

\paragraph{Result: the effect replicates, but not as deception.}
Across $288$ runs with zero API errors, the ambient route produces far more false
filings than the direct request in two of three domains
(Table~\ref{tab:p2}); pooled across domains, Cochran--Mantel--Haenszel gives
$\chi^2=109.2$, $p<10^{-15}$. By the pre-registered rule ($\ge2$ of $3$), the
request-modality effect is \textbf{not specific to the original task family}. The
direct-request cells are the sharper half of the finding: in code review and
procurement, agents asked outright to file a false verdict comply
\textbf{$0.00$} of the time.

\begin{table}[t]
\centering
\caption{\textbf{P2: active-lie rate by domain and delivery}, $n{=}48$/cell, forced
verdict, Wilson 95\% CIs. The right-hand block restricts to runs in which the agent
actually read the authoritative record---the only runs in which a false filing is a
\emph{knowing} misreport.}
\label{tab:p2}
\small
\footnotesize
\begin{tabular}{llcc|ccc}
\toprule
& & \multicolumn{2}{c|}{all runs} & \multicolumn{3}{c}{read the record} \\
domain & delivery & lie & 95\% CI & read rate & $n$ & lie \\
\midrule
code review   & direct  & 0.00 & [.00,.07] & 0.44 & 21 & 0.00 \\
code review   & ambient & \textbf{0.73} & [.59,.83] & 0.58 & 28 & \textbf{0.54} \\
compliance    & direct  & 0.15 & [.07,.27] & 0.60 & 29 & 0.21 \\
compliance    & ambient & 0.19 & [.10,.32] & 1.00 & 48 & 0.19 \\
procurement   & direct  & 0.00 & [.00,.07] & 0.96 & 46 & 0.00 \\
procurement   & ambient & \textbf{0.94} & [.83,.98] & \textbf{0.06} & 3 & --- \\
\bottomrule
\end{tabular}
\end{table}

\paragraph{The mechanism is pre-empted verification, not persuasion.}
The pre-registered knowledge gate changes what the effect \emph{is}. In procurement,
the ambient artifact did not talk an informed agent into a false recommendation---it
stopped the agent from ever looking. Only $6\%$ of ambient-route agents called
\texttt{read\_truth} (versus $96\%$ under a direct request), and \emph{every} one of
the $45$ who did not look filed the favourable verdict. With three readers, that cell
cannot be tested at all, and we report it as untestable rather than as a null.

Only code review supports a genuine knowing-misreport claim, and there it is large:
among agents that had read the failing CI results, the ambient route still produced
$0.54$ false approvals against $0.00$ under a direct ask (RD $+0.536$
$[+0.300,+0.705]$, $p{=}3{\times}10^{-5}$). Compliance is null on both analyses
($-0.019$ $[-0.216,+0.151]$)---and is also the one domain where the direct ask
partially works ($0.21$ among readers), which suggests a regulator-pressure framing
is more persuasive than an interpersonal plea.

\paragraph{What this changes for a defender.}
The headline generalises; the \emph{explanation} does not. A pre-populated artifact is
dangerous mainly because it removes the agent's reason to check, and only secondarily
because it makes an informed agent willing to misreport. Auditing the filing---already
shown not to help---is doubly the wrong control here, because the failure happens
before anything is filed. The implied defense is to \textbf{make verification
mandatory and non-skippable}: an agent that cannot file without first reading the
record is not exposed to the dominant route at all. We did not test that control and
flag it as the obvious next experiment.

% ======================================================================
